\documentclass[12pt,a4paper]{article}
\usepackage[utf8]{inputenc}
\usepackage[portuguese]{babel}
\usepackage{amsmath,amsfonts,amssymb}
\usepackage{graphicx}
\usepackage{float}
\usepackage{booktabs}
\usepackage{multirow}
\usepackage{array}
\usepackage{geometry}
\usepackage{fancyhdr}
\usepackage{algorithm}
\usepackage{algorithmic}
\usepackage{listings}
\usepackage{xcolor}

\geometry{margin=2.5cm}
\pagestyle{fancy}
\fancyhf{}
\rhead{\thepage}
\lhead{Classificadores Baseados em Clustering}

% Configuração para código
\lstset{
    language=Python,
    basicstyle=\ttfamily\small,
    keywordstyle=\color{blue},
    commentstyle=\color{green},
    stringstyle=\color{red},
    numbers=left,
    numberstyle=\tiny,
    frame=single,
    breaklines=true
}

\title{
    \textbf{Classificadores Multiclasse Baseados em\\Algoritmos de Agrupamento}\\
    \vspace{0.5cm}
    \large Inteligência Computacional
}

\author{
    [Seu Nome]\\
    [Sua Matrícula]\\
    CEFET-MG
}

\date{\today}

\begin{document}

\maketitle

\begin{abstract}
Este trabalho apresenta a implementação e avaliação de três classificadores multiclasse baseados em algoritmos de agrupamento com treinamento supervisionado. Os algoritmos implementados foram K-means, Gaussian Mixture Model (GMM) e Clustering Hierárquico. O método do cotovelo foi utilizado para determinar o número ideal de clusters. Os classificadores foram avaliados nos datasets Adult e Dry Bean do repositório UCI, com 30 repetições cada, utilizando divisão 70\%-30\% para treino-teste. Os resultados mostram que [completar com resultados obtidos].
\end{abstract}

\section{Introdução}

Os algoritmos de agrupamento (clustering) são tradicionalmente utilizados para análise exploratória de dados e descoberta de padrões não supervisionados. No entanto, quando combinados com informações de classe durante o treinamento, podem ser adaptados para tarefas de classificação supervisionada.

Este trabalho implementa três classificadores baseados em clustering:
\begin{itemize}
    \item \textbf{K-means}: Algoritmo de particionamento baseado em centroides
    \item \textbf{Gaussian Mixture Model (GMM)}: Modelo probabilístico baseado em misturas gaussianas
    \item \textbf{Clustering Hierárquico}: Algoritmo aglomerativo baseado em distâncias
\end{itemize}

\section{Metodologia}

\subsection{Pseudocódigo do Classificador}

\begin{algorithm}[H]
\caption{Classificador Baseado em Clustering}
\begin{algorithmic}[1]
\REQUIRE $X_{treino}$, $y_{treino}$, $X_{teste}$, algoritmo
\ENSURE $y_{pred}$
\STATE Normalizar $X_{treino}$ usando StandardScaler
\STATE Encontrar $k_{ótimo}$ usando método do cotovelo:
\FOR{$k = 2$ \TO $k_{max}$}
    \STATE Aplicar clustering com $k$ clusters
    \STATE Calcular inércia/score
\ENDFOR
\STATE Encontrar cotovelo na curva de inércia
\STATE Aplicar clustering com $k_{ótimo}$ em $X_{treino}$
\FOR{cada cluster $i$}
    \STATE Encontrar classe mais frequente no cluster
    \STATE Mapear $cluster_i \rightarrow classe_{frequente}$
\ENDFOR
\STATE Normalizar $X_{teste}$
\FOR{cada ponto em $X_{teste}$}
    \STATE Atribuir ao cluster mais próximo
    \STATE Mapear cluster $\rightarrow$ classe
\ENDFOR
\RETURN $y_{pred}$
\end{algorithmic}
\end{algorithm}

\subsection{Método do Cotovelo}

O método do cotovelo é utilizado para determinar automaticamente o número ideal de clusters. O algoritmo calcula a inércia (soma das distâncias quadráticas aos centroides) para diferentes valores de $k$ e identifica o ponto onde a taxa de redução da inércia diminui significativamente, formando um "cotovelo" na curva.

\subsection{Treinamento Supervisionado}

Após o agrupamento, cada cluster é associado à classe mais frequente entre os pontos que o compõem. Esta abordagem permite que algoritmos não supervisionados sejam utilizados para classificação.

\subsection{Datasets Utilizados}

\subsubsection{Adult Dataset}
\begin{itemize}
    \item \textbf{Amostras}: 48,842
    \item \textbf{Features}: 14
    \item \textbf{Classes}: 2 (<=50K, >50K)
    \item \textbf{Tipo}: Classificação binária de renda
\end{itemize}

\subsubsection{Dry Bean Dataset}
\begin{itemize}
    \item \textbf{Amostras}: 13,611
    \item \textbf{Features}: 16
    \item \textbf{Classes}: 7 tipos de feijão
    \item \textbf{Tipo}: Classificação multiclasse
\end{itemize}

\subsection{Protocolo Experimental}

\begin{itemize}
    \item \textbf{Repetições}: 30 execuções com sementes 1-30
    \item \textbf{Divisão}: 70\% treino, 30\% teste
    \item \textbf{Normalização}: StandardScaler aplicado às features
    \item \textbf{Métricas}: Acurácia, desvio padrão, matriz de confusão
\end{itemize}

\section{Resultados Experimentais}

\subsection{Dataset Adult}

\begin{table}[H]
\centering
\caption{Resultados para o Dataset Adult}
\begin{tabular}{lcccc}
\toprule
\textbf{Algoritmo} & \textbf{Acurácia Média} & \textbf{Desvio Padrão} & \textbf{Min} & \textbf{Max} \\
\midrule
K-means & 0.XXXX & 0.XXXX & 0.XXXX & 0.XXXX \\
GMM & 0.XXXX & 0.XXXX & 0.XXXX & 0.XXXX \\
Hierárquico & 0.XXXX & 0.XXXX & 0.XXXX & 0.XXXX \\
\bottomrule
\end{tabular}
\end{table}

\begin{figure}[H]
\centering
\includegraphics[width=0.8\textwidth]{../img/adult_comprehensive_results.png}
\caption{Resultados detalhados para o dataset Adult}
\end{figure}

\subsection{Dataset Dry Bean}

\begin{table}[H]
\centering
\caption{Resultados para o Dataset Dry Bean}
\begin{tabular}{lcccc}
\toprule
\textbf{Algoritmo} & \textbf{Acurácia Média} & \textbf{Desvio Padrão} & \textbf{Min} & \textbf{Max} \\
\midrule
K-means & 0.XXXX & 0.XXXX & 0.XXXX & 0.XXXX \\
GMM & 0.XXXX & 0.XXXX & 0.XXXX & 0.XXXX \\
Hierárquico & 0.XXXX & 0.XXXX & 0.XXXX & 0.XXXX \\
\bottomrule
\end{tabular}
\end{table}

\begin{figure}[H]
\centering
\includegraphics[width=0.8\textwidth]{../img/dry_bean_comprehensive_results.png}
\caption{Resultados detalhados para o dataset Dry Bean}
\end{figure}

\subsection{Método do Cotovelo}

\begin{figure}[H]
\centering
\includegraphics[width=0.6\textwidth]{../img/elbow_method_example.png}
\caption{Exemplo do método do cotovelo para determinação do número ideal de clusters}
\end{figure}

\subsection{Matrizes de Confusão}

\begin{figure}[H]
\centering
\begin{minipage}{0.45\textwidth}
\centering
\includegraphics[width=\textwidth]{../img/confusion_matrix_adult.png}
\caption{Matriz de confusão média - Adult}
\end{minipage}
\hfill
\begin{minipage}{0.45\textwidth}
\centering
\includegraphics[width=\textwidth]{../img/confusion_matrix_bean.png}
\caption{Matriz de confusão média - Dry Bean}
\end{minipage}
\end{figure}

\section{Análise dos Resultados}

\subsection{Desempenho dos Algoritmos}

\textbf{K-means}: [Análise do desempenho do K-means]

\textbf{GMM}: [Análise do desempenho do GMM]

\textbf{Clustering Hierárquico}: [Análise do desempenho do clustering hierárquico]

\subsection{Comparação entre Datasets}

O dataset Adult, sendo um problema de classificação binária, apresentou [análise]. Já o dataset Dry Bean, com 7 classes, mostrou [análise].

\subsection{Eficácia do Método do Cotovelo}

O método do cotovelo mostrou-se eficaz para [análise da eficácia].

\subsection{Estabilidade dos Resultados}

A análise dos desvios padrão ao longo das 30 repetições indica [análise da estabilidade].

\section{Conclusões}

Este trabalho demonstrou que classificadores baseados em algoritmos de clustering podem ser eficazes quando combinados com treinamento supervisionado. Os principais achados incluem:

\begin{itemize}
    \item O método do cotovelo é uma ferramenta útil para determinação automática do número de clusters
    \item A normalização dos dados é crucial para o desempenho dos algoritmos
    \item [Outras conclusões baseadas nos resultados]
\end{itemize}

\subsection{Limitações}

\begin{itemize}
    \item Dependência da qualidade do clustering inicial
    \item Sensibilidade à normalização dos dados
    \item Complexidade computacional do clustering hierárquico
    \item Suposições de distribuição gaussiana no GMM
\end{itemize}

\subsection{Trabalhos Futuros}

\begin{itemize}
    \item Investigação de outros algoritmos de clustering
    \item Otimização de hiperparâmetros
    \item Aplicação em datasets com mais classes
    \item Comparação com classificadores tradicionais
\end{itemize}

\section{Referências}

\begin{thebibliography}{9}

\bibitem{kmeans}
MacQueen, J. (1967). Some methods for classification and analysis of multivariate observations. In Proceedings of the fifth Berkeley symposium on mathematical statistics and probability (Vol. 1, No. 14, pp. 281-297).

\bibitem{gmm}
Reynolds, D. A. (2009). Gaussian mixture models. Encyclopedia of biometrics, 741, 659-663.

\bibitem{hierarchical}
Murtagh, F., \& Contreras, P. (2012). Algorithms for hierarchical clustering: an overview. Wiley Interdisciplinary Reviews: Data Mining and Knowledge Discovery, 2(1), 86-97.

\bibitem{elbow}
Thorndike, R. L. (1953). Who belongs in the family?. Psychometrika, 18(4), 267-276.

\bibitem{uci}
Dua, D. and Graff, C. (2019). UCI Machine Learning Repository. Irvine, CA: University of California, School of Information and Computer Science.

\end{thebibliography}

\end{document}